% !PRINTED none

\begin{center}
\gesperrt{First Chapter.}

\vspace{1em}
\gesperrt{General Group Theory.}

\rule{2cm}{0.4pt}

\vspace{1em}
\S.\,1.

\gesperrt{Definition of Groups.}
\end{center}

\noindent In the first volume, in connection with the permutations, we became acquainted with the concept of a group, and we made important algebraic applications of it. It must now be our next task to grasp this concept, so exceedingly important in the whole of recent mathematics, in a more general way, and to become acquainted with the laws that govern it. We place the following definition at the head:

A system $P$ of things (elements) of any kind whatever becomes a group when the following conditions are fulfilled:

\noindent 1.\ \gesperrt{A rule is given, by which from a first and a second element of the system a wholly determinate third element of the same system is derived.}

One writes symbolically, if $a$ is the first, $b$ the second, $c$ the third element:
\[
ab = c, \quad c = ab,
\]
and calls $c$ \gesperrt{compounded} from $a$ and $b$, and $a$ and $b$ the \gesperrt{components} of $c$.

In this composition the commutative law is not in general presupposed, that is, $ab$ may be different from $ba$; on the other hand,

\noindent 2.\ \gesperrt{the associative law is presupposed,} that is, if $a, b, c$ are any three elements of $P$, then
\[
(ab)c = a(bc),
\]
%CONT>>
